% ---
% RESUMOS
% ---

% resumo em português
\setlength{\absparsep}{18pt} % ajusta o espaçamento dos parágrafos do resumo
\begin{resumo}
	A qualidade da educação básica é frequentemente avaliada de forma isolada, sem considerar os fatores socioeconômicos e ambientais que impactam o desempenho escolar fora do ambiente de sala de aula. Neste contexto, os indicadores do Índice de Desenvolvimento Sustentável das Cidades (IDSC) oferecem uma perspectiva multidimensional que pode ser explorada para compreender o desempenho educacional em nível municipal. Este trabalho propõe analisar o desempenho de modelos preditivos baseados em árvores de decisão na tarefa de classificação das faixas de desempenho do Índice de Desenvolvimento da Educação Básica (IDEB) nos municípios do estado de Santa Catarina, com dados do IDSC. A metodologia compreende etapas de análise exploratória, pré-processamento com remoção preventiva de vazamento de dados, seleção de atributos por filtragem de correlação e remoção dos indicadores com relação à educação, além de validação cruzada para ajuste de hiperparâmetros. Foram comparados os algoritmos \textit{Random Forest}, \textit{XGBoost}, \textit{LightGBM} e \textit{CatBoost}, avaliados separadamente para os Anos Iniciais (AI) e para os Anos Finais (AF) do ensino fundamental, por meio de Acurácia, Precisão, \textit{Recall} e F1-\textit{score} (macro). Algoritmos do tipo \textit{Gradient Boosting}, obtiveram o melhor desempenho preditivo nas duas variáveis alvo, embora com diferenças pequenas entre os modelos. A análise de importância dos atributos mostrou que os indicadores de saúde pública (ODS 3) concentram a maior parcela da importância na separação das faixas do IDEB nos Anos Finais. Portanto, os resultados indicam associação entre o contexto socioeconômico dos municípios catarinenses e as faixas de desempenho do IDEB, fornecendo embasamentos quantitativos e interpretáveis para o planejamento e a formulação de políticas públicas integradas na gestão educacional.
	
	\textbf{Palavras-chave}: Índice de Desenvolvimento da Educação Básica; Aprendizado de Máquina; Objetivos de Desenvolvimento Sustentável; Interpretabilidade de Modelos; Políticas Públicas Educacionais.
\end{resumo}

